% ECONOMICS_PROBLEM: {"id": "I13-462", "title": "Long-term-care finance and informal care", "jel_code": "I13", "source_id": "OP-0059"}
% !TeX program = xelatex
\documentclass[11pt,a4paper]{article}
\usepackage[margin=23mm]{geometry}
\usepackage{fontspec}
\setmainfont{Latin Modern Roman}
\usepackage{amsmath,amssymb,mathtools}
\usepackage{xcolor,enumitem,booktabs,longtable,array,xurl,hyperref}
\definecolor{muted}{HTML}{53636C}
\hypersetup{colorlinks=true,urlcolor=blue,linkcolor=blue}
\setlength{\parindent}{0pt}
\setlength{\parskip}{5pt}
\newcommand{\R}{\mathbb R}
\newcommand{\E}{\mathbb E}
\newcommand{\N}{\mathbb N}
\newcommand{\ind}{\mathbf 1}
\newcommand{\Var}{\operatorname{Var}}
\newcommand{\Cov}{\operatorname{Cov}}
\newcommand{\supp}{\operatorname{supp}}
\DeclareMathOperator*{\argmin}{arg\,min}
\DeclareMathOperator*{\argmax}{arg\,max}
\newcommand{\entrymeta}[3]{{\sffamily\small\color{muted}\textbf{\texttt{#1}}\quad|\quad #2\par #3\par}}
\newcommand{\Problemref}[1]{\texttt{#1}}
\newcommand{\JELref}[2]{\texttt{#2} (JEL \texttt{#1})}
\begin{document}
\section*{Long-term-care finance and informal care}
\textbf{Problem ID:} \texttt{I13-462}\quad \textbf{JEL:} \texttt{I13}\par
% BEGIN ECONOMICS STATEMENT
\label{problem:OP-0059}
\label{atlas:prob:H9}
\noindent\textbf{Original atlas identifier:} H9; entry 59. \textbf{Field:} Health economics. \textbf{Original tractability label:} Hard.

\noindent\textbf{Classification:} Parameterized theoretical/empirical research agenda; not a certified unsolved theorem.

\subsection*{Definitions and assumptions}
A family consists of an older adult and a specified set of potential caregivers. At date $t$, state $s_t$ includes disability, cognition, assets, insurance coverage, caregiver health, and wages. Actions comprise formal-care quantity $f_t$, informal-care hours $h_{jt}$ for caregiver $j$, labor supply, saving, and insurance purchase. A dynamic model $m$ specifies care production, health transitions, bargaining/altruism within the family, private insurer selection and premiums, and equilibrium selection. A public policy $p$ specifies eligibility, copayments, cash or in-kind benefits, caregiver compensation, and financing. Let $u_a$, $u_j$ be cardinal adult/caregiver utilities, $w_a,w_j\geq0$ fixed social weights, $\beta\in(0,1)$ discount factor, and $G_t(p;m)$ net public outlays; $\lambda$ is their utility-unit social cost.

\subsection*{Formal research question}
\emph{Editorial formulation: the mathematical target below is not a theorem quoted from the references.}

For legal policies obeying a stated present-value financing constraint and verifiable-need rules, characterize
\[
 \sup_{p\in\mathcal P}\mathbb E_m\sum_{t=0}^{\infty}\beta^t
 \left[w_a u_a(s_t,f_t,\mathbf h_t)+\sum_jw_j u_j(s_t,h_{jt})-\lambda G_t(p;m)\right].
\]
Utilities include consumption and foregone leisure/earnings consistently; do not subtract caregiver time costs again. Determine whether reforms separately identify substitution/complementarity between formal and informal care, insurer crowd-out, and caregiver labor responses. Compare cash-benefit, in-kind, and private-insurance designs over models consistent with family-linked data. If family bargaining or latent need is not point identified, report bounds on welfare and $\varepsilon$-optimal policies rather than a single universal division of risk.

\subsection*{Interpretation and scope}
Risk sharing is multidimensional: financing a nursing-home bill, insuring caregiver earnings, and guaranteeing access under cognitive impairment are different policy functions. Informal care is a real input with opportunity costs, not a free fiscal saving. Family-member transfers cancel in aggregate resource accounting but matter under unequal welfare weights and bargaining. Formal and informal care can be complements for some tasks and substitutes for others; a single average elasticity may miss this distinction. Public coverage can change private insurance demand through eligibility rules and family options, so financing comparisons need all three institutions in the same model. Cross-country institutional descriptions alone do not identify reform effects. The extension asks for joint family/insurance/public-budget policy evaluation with verified needs and caregiver welfare.

\subsection*{Source audit and status}
[S1] and [S2] are verified; [S2] is the November 2023 NBER working paper by Gruber, McGarry, and Hanzel. The exact Mommaerts (2018) citation cannot be verified from the incomplete original. [S4] supplies a separately identified related 2025 publication; its November 2024 online date is distinguished from its January 2025 issue.

\subsection*{References}
\begin{enumerate}[label={[S\arabic*]},leftmargin=*,itemsep=0.6em]
\item Jeffrey R. Brown and Amy Finkelstein (2008). The Interaction of Public and Private Insurance: Medicaid and the Long-Term Care Insurance Market. American Economic Review 98(3), 1083--1102. DOI: 10.1257/aer.98.3.1083.
\url{https://www.aeaweb.org/articles?id=10.1257/aer.98.3.1083}
\newline\emph{Locator and support limits:} Abstract and model: Medicaid and private long-term-care insurance interaction; family caregiving is an additional margin.
\item Jonathan Gruber, Kathleen M. McGarry, and Charles Hanzel (2023). Long-Term Care Around the World. NBER Working Paper 31882, November. DOI: 10.3386/w31882.
\url{https://www.nber.org/papers/w31882}
\newline\emph{Locator and support limits:} Abstract and report: cross-country institutional comparisons; not causal evidence that one financing rule dominates.
\item Unresolved original: Mommaerts (2018). No title, venue, or identifier is supplied by the atlas; an exact 2018 version was not verified.
\url{https://conference.nber.org/confer/2017/CAHs17/Mommaerts.pdf}
\newline\emph{Locator and support limits:} A 2017 conference manuscript and a 2025 publication with the relevant title were found; neither proves an exact 2018 version.
\item Corina Mommaerts (2025). Long-Term Care Insurance and the Family. Journal of Political Economy 133(1), 1--52. DOI: 10.1086/732887.
\url{https://www.journals.uchicago.edu/doi/10.1086/732887}
\newline\emph{Locator and support limits:} Abstract and article: family care and insurance in a dynamic parent-child model. Online 7 November 2024; issue January 2025.
\end{enumerate}

\subsection*{Common conventions: Source mathematical conventions}

Unless an entry states otherwise, agent and firm sets are finite. State and action spaces are measurable, strategies and outcome maps are measurable, probability laws are specified, and expectations exist for the targets under discussion. Infinite-horizon discount factors lie in $(0,1)$ and discounted objectives are finite. Norms, tolerances and complexity input sizes must be specified for computational comparisons. Constants or primitives that appear locally are inputs, not empirical facts asserted by this catalogue.

\textbf{Identification.} A statistical model $m\in\mathcal M$ implies an observable law $P_m$ and target $\theta(m)$. At law $P$, its identified set is
\[
\Theta(P;\mathcal M)=\{\theta(m):m\in\mathcal M,\ P_m=P\}.
\]
Point identification holds when this set is a singleton, or equivalently when $P_{m_1}=P_{m_2}$ implies $\theta(m_1)=\theta(m_2)$ for all compatible models. Sharp bounds mean bounds attained, or approached when the boundary is not attained, by compatible models. Writing this set is a definition, not a solution: the substantive task is to establish informative restrictions and derive or estimate the set. An unrestricted-confounding model may give only trivial bounds.

\textbf{Inference.} Identification, estimability and valid uncertainty quantification are separate questions. Fix $\alpha\in(0,1)$. A confidence procedure $C_n$ over a declared law class $\mathcal P$ has asymptotic uniform coverage at level $1-\alpha$ when
\[
\liminf_{n\to\infty}\inf_{P\in\mathcal P}\Pr_P\{\theta(P)\in C_n\}\geq1-\alpha.
\]
For set-identified targets the coverage event must be defined separately, for example coverage of the full identified set. If an entry uses identification-indexed models instead of uniquely indexed laws, the same coverage convention is applied over those models. Pointwise limits do not establish uniform validity, and asymptotic validity does not guarantee finite-sample precision. Sample sizes, dependence structures and weak-identification sequences are part of each application.

\textbf{Counterfactuals.} $Y_i(a)$ is a potential outcome under a specified intervention, including its timing, implementation and versions. It is not $Y_i$ conditional on voluntarily choosing $a$. A full intervention vector permits interference; shorthand scalar treatment notation does not silently impose no interference. Assignment, selection, attrition, anticipation and measurement must be specified. A claim about an instrument's local effect is not automatically a population policy effect.

\textbf{Economic equilibrium.} A counterfactual model includes best responses, constraints, feasibility, market clearing, state transitions and expectations consistent with the stated information. When several equilibria exist, either a declared selector is an input or the target is set-valued. Comparisons across policy regimes must use the stated selection convention. Reduced-form relationships are not presumed invariant after an equilibrium-changing intervention.

\textbf{Optimization and welfare.} Optimization is over the stated feasible set. If attainment is not established, $\sup$ or $\inf$ and $\varepsilon$-optimal choices replace an asserted maximizer or minimizer. For example, with value $V=\sup_{a\in\mathcal A}W(a)<\infty$, an $\varepsilon$-optimal set is $\{a\in\mathcal A:W(a)\geq V-\varepsilon\}$ for $\varepsilon>0$. Benefits, costs, risk and health or environmental outcomes can be aggregated only using declared units and conversion weights. Interpersonal utility weights, the treatment of fiscal transfers and foreign profits, discounting, and the behavioral welfare criterion are explicit inputs. A comparative-static derivative additionally requires existence and appropriate differentiability of the selected counterfactual.

\textbf{Sources.} Local labels [S1], [S2], and so forth restart in every question. Locators identify the relevant abstract, model, section or result at the level actually checked; a bibliographic or abstract-level verification is not represented as inspection of an entire proof. Working papers are distinguished from later publications and changing versions. Source summaries are paraphrased. All URL checks and editorial formulations refer to the audit date stated above.
% END ECONOMICS STATEMENT
\end{document}
